\documentclass{szb-solution}

\area{Komplex számok}
\title{Komplex számok}
\subject{Matematika G1}
\subjectCode{BMETE94BG01}
\date{Utoljára frissítve: \today}
\docno{3}

\begin{document}
\maketitle

\subsection{Komplex kifejezések egyszerűsítése}

\begin{enumerate}[a)]
  \item Mivel $e^{\iu\pi/3}=\frac12+\frac{\sqrt3}{2}\iu$, ezért
        \begin{align*}
          \overline{\left(\frac{2-\iu}{e^{\iu\pi/3}}\right)}
          &=\frac{2+\iu}{e^{-\iu\pi/3}}
          =(2+\iu)e^{\iu\pi/3}
          \\
          &=(2+\iu)\left(\frac12+\frac{\sqrt3}{2}\iu\right)
          =\boxed{1-\frac{\sqrt3}{2}
          +\left(\sqrt3+\frac12\right)\iu}
          \text.
        \end{align*}

  \item Először
        $$
          \frac{5+\iu}{3-2\iu}
          =\frac{(5+\iu)(3+2\iu)}{13}
          =1+\iu
          =\sqrt2e^{\iu\pi/4}
          \text.
        $$
        A konjugált tényező:
        $$
          \overline{
            3e^{\iu\pi/4}\cdot2e^{\iu3\pi/2}\cdot e^{\iu5\pi/12}
          }
          =\overline{6e^{\iu13\pi/6}}
          =6e^{-\iu\pi/6}
          \text.
        $$
        Így a teljes kifejezés
        $$
          6\sqrt2e^{\iu(\pi/4-\pi/6)}
          =6\sqrt2e^{\iu\pi/12}
          =\boxed{3(\sqrt3+1)+3(\sqrt3-1)\iu}
          \text.
        $$

  \item A tényezők trigonometrikus alakban:
        $$
          \frac{5e^{\iu7\pi/13}}{4e^{\iu3\pi/4}},
          \qquad
          \overline{\left(\frac{1}{2\iu}\right)}=\frac{\iu}{2}
          =\frac12e^{\iu\pi/2},
          \qquad
          2\sqrt3+2\iu=4e^{\iu\pi/6}
          \text.
        $$
        Ezért
        \begin{align*}
          \frac{5e^{\iu7\pi/13}}{4e^{\iu3\pi/4}}
          \cdot\frac12e^{\iu\pi/2}\cdot4e^{\iu\pi/6}
          &=\frac52e^{\iu(7\pi/13-3\pi/4+\pi/2+\pi/6)}
          \\
          &=\boxed{\frac52e^{\iu71\pi/156}}
          \text.
        \end{align*}
\end{enumerate}

\subsection{Hatványozás}

\begin{enumerate}[a)]
  \item Mivel $\iu-1=\sqrt2e^{\iu3\pi/4}$, a Moivre-formula alapján
        $$
          (\iu-1)^{16}
          =(\sqrt2)^{16}e^{\iu12\pi}
          =\boxed{256}
          \text.
        $$

  \item Vegyük észre, hogy $21-35\iu=7(3-5\iu)$ és
        $$
          \frac{1+\iu}{1-\iu}=\iu,
          \qquad \iu^4=1
          \text.
        $$
        Ezért
        \begin{align*}
          &(3+5\iu)^4(21-35\iu)^5
          \left(\frac{1+\iu}{1-\iu}\right)^4
          \\
          &\qquad
          =7^5(3+5\iu)^4(3-5\iu)^5
          =7^5\bigl((3+5\iu)(3-5\iu)\bigr)^4(3-5\iu)
          \\
          &\qquad
          =\boxed{7^5\cdot34^4(3-5\iu)}
          \text.
        \end{align*}
\end{enumerate}

\subsection{Gyökvonás}

\begin{enumerate}[a)]
  \item A $-8=8e^{\iu(\pi+2k\pi)}$ harmadik gyökei
        $$
          z_k=2e^{\iu(\pi+2k\pi)/3},
          \qquad k=0,1,2
          \text,
        $$
        vagyis
        $$
          \boxed{\sqrt[3]{-8}=\{1+\sqrt3\iu;\,-2;\,1-\sqrt3\iu\}}
          \text.
        $$

  \item Az egység negyedik gyökei
        $$
          z_k=e^{\iu k\pi/2},
          \qquad k=0,1,2,3
          \text,
        $$
        tehát
        $$
          \boxed{\sqrt[4]{1}=\{1;\,\iu;\,-1;\,-\iu\}}
          \text.
        $$

  \item Legyen $(x+y\iu)^2=3+4\iu$. Ekkor
        $$
          x^2-y^2=3,
          \qquad 2xy=4
          \text.
        $$
        A rendszer valós megoldásai $(x,y)=(2,1)$ és $(-2,-1)$, ezért
        $$
          \boxed{\sqrt{3+4\iu}=\{2+\iu;\,-2-\iu\}}
          \text.
        $$
\end{enumerate}

\subsection{Komplex egyenletek}

\begin{enumerate}[a)]
  \item A $z^4=81\iu=81e^{\iu(\pi/2+2k\pi)}$ egyenlet gyökei
        $$
          \boxed{
            z_k=3e^{\iu(\pi/8+k\pi/2)},
            \qquad k=0,1,2,3
          }
          \text.
        $$

  \item A másodfokú megoldóképlettel
        $$
          z_{1,2}
          =\frac{6\pm\sqrt{36-52}}{2}
          =\frac{6\pm4\iu}{2}
          =\boxed{3\pm2\iu}
          \text.
        $$
\end{enumerate}

\subsection{Komplex egyenletrendszer}

Az első egyenletből $z_1=1+\iu-2z_2$. Ezt a másodikba helyettesítve
$$
  3(1+\iu-2z_2)+\iu z_2=2-3\iu,
$$
$$
  (-6+\iu)z_2=-1-6\iu
  \quad\Longrightarrow\quad
  z_2=\frac{-1-6\iu}{-6+\iu}=\iu
  \text.
$$
Végül $z_1=1+\iu-2\iu=1-\iu$, tehát
$$
  \boxed{z_1=1-\iu,
  \qquad z_2=\iu}
  \text.
$$

\subsection{Geometriai helyek}

Írjuk $z=x+y\iu$ alakban.

\begin{enumerate}[a)]
  \item Az $\iIm\{z\}>0$ feltétel a komplex számsík nyílt felső félsíkját,
        vagyis az $y>0$ pontokat adja.

  \item Az $|z-a|=|z-b|$ feltétel az $a$ és $b$ pontoktól egyenlő távolságra
        levő pontokat írja le: ez az $ab$ szakasz felezőmerőlegese
        ($a\neq b$ esetén).

  \item Az $|z|<1-\iRe\{z\}$ feltétel
        $$
          \sqrt{x^2+y^2}<1-x
        $$
        alakú. Négyzetre emelve
        $$
          x^2+y^2<(1-x)^2
          \quad\Longleftrightarrow\quad
          \boxed{x<\frac{1-y^2}{2}}
          \text.
        $$
        Ez az $x=(1-y^2)/2$ balra nyíló parabola bal oldalán fekvő nyílt
        tartomány; a határparabola nem része a halmaznak.
\end{enumerate}

\subsection{Négyzet további csúcsai}

A két szomszédos csúcs $z_1=3+2\iu$ és $z_2=5+4\iu$, így az oldal vektora
$$
  d=z_2-z_1=2+2\iu
  \text.
$$
A rá merőleges, azonos hosszúságú oldal $\pm\iu d$. Az egyik irányban
$$
  \iu d=-2+2\iu,
  \qquad
  z_1+\iu d=1+4\iu,
  \qquad
  z_2+\iu d=3+6\iu
  \text,
$$
a másik irányban
$$
  -\iu d=2-2\iu,
  \qquad
  z_1-\iu d=5,
  \qquad
  z_2-\iu d=7+2\iu
  \text.
$$
Tehát a további csúcsok
$$
  \boxed{\{1+4\iu,\,3+6\iu\}}
  \quad\text{vagy}\quad
  \boxed{\{5,\,7+2\iu\}}
  \text.
$$

\subsection{Kör a komplex számsíkon}

A $(-2;1)$ középponthoz a $-2+\iu$ komplex szám tartozik. A $4$ sugarú kör
egyenlete
$$
  \boxed{|z-(-2+\iu)|=4}
  \quad\Longleftrightarrow\quad
  \boxed{|z+2-\iu|=4}
  \text.
$$
Ha $z=x+y\iu$, akkor ugyanez
$$
  \boxed{(x+2)^2+(y-1)^2=16}
  \text.
$$

\end{document}
